\section{Discusión}

La optimización comprehensiva de hiperparámetros conducida vía el framework Optuna, en múltiples conjuntos de datos dispares, proporciona una confirmación robusta de la hipótesis de trabajo. Los resultados sugieren un cambio de paradigma en interpolación EEG: métodos conscientes de grafos y de regularización temporal (TV) no son meramente competitivos teóricamente, sino prácticamente superiores en escenarios realistas de canales faltantes cuando están apropiadamente sintonizados.

\subsection{Fidelidad Fisiológica y Temporal (ERP \& PSD)}
Un hallazgo crítico de la fase de optimización es el impacto profundo de regularización temporal en características fisiológicas de señal. Incluso cuando líneas base geométricas (p. ej., Spherical Spline o RBFI) son exhaustivamente optimizadas para minimizar su error espacial, funcionan fundamentalmente como suavizadores espaciales. En contraste, TRSS incorpora conocimiento previo de la suavidad temporal inherente en generadores neuro-eléctricos.

Esta ventaja es más evidente en la preservación de Potenciales Relacionados con Eventos (ERPs) y Densidad Espectral de Potencia (PSD). Como fue mostrado previamente en la sección de Resultados (\figurename~\ref{fig:erp_optimized} y \figurename~\ref{fig:psd_optimized}), TRSS recupera con precisión amplitud y latencia de componentes ERP (como P300 o N100) incluso bajo pérdida severa de canales contiguos (40\% de canales cercanos faltantes). Líneas base optimizadas para mejores casos espaciales sufren distorsión de fase y atenuación de amplitud en el dominio temporal, comprometiendo tareas posteriores clínicas o BCI.

\subsection{Generalización Multi-Dataset}
La varianza en desempeño en datasets subraya la importancia del framework de optimización Optuna. En configuraciones densamente muestreadas y altamente correlacionadas (PhysioNet, 64 canales), los modelos espaciales funcionan adecuadamente. Sin embargo, en setups de densidad más baja o aquellos con dinámicas temporales complejas (BCI IV 2a y MNE Sample), los priors temporales de TRSS proporcionan un buffer sustancial contra fallo catastrófico, particularmente cuando clusters cercanos de electrodos se pierden.

\subsection{Posicionamiento Contra Enfoques de Aprendizaje Profundo}
Es importante contextualizar estos resultados contra el panorama más amplio de métodos de reconstrucción EEG. Enfoques recientes de aprendizaje profundo---incluyendo autoencoders convolucionales, redes neuronales de grafos (GNNs), y arquitecturas basadas en transformers---han mostrado promesa en tareas de imputación de señales relacionadas. Sin embargo, una comparación directa fue intencionalmente excluida de este benchmark por dos razones principiadas. Primera, métodos GSP como TRSS son \emph{sin entrenamiento}: no requieren fase de pre-entrenamiento, sin datos etiquetados, y sin calibración específica de sujeto, haciéndolos inmediatamente desplegables en cualquier montaje EEG. Segunda, la convergencia teórica del descenso de gradiente TRSS está garantizada por el límite de Lipschitz (Eq.~\ref{eq:lipschitz}), proporcionando interpretabilidad y confiabilidad que redes neuronales caja-negra no pueden ofrecer. Trabajo futuro debería investigar arquitecturas híbridas que combinen los sesgos inductivos de GSP con el poder representacional de modelos aprendidos.

\subsection{Limitaciones}
Este estudio tiene limitaciones que deben ser reconocidas. Primera, la simulación de canales faltantes usa enmascaramiento binario (canal totalmente presente o ausente), mientras que degradación real es frecuentemente gradual. Segunda, todos los datasets corresponden a sujetos adultos sanos; validación en poblaciones clínicas (epilepsia, neonatos) permanece como pregunta abierta. Tercera, la evaluación actual usa una única división temporal por dataset; trabajo futuro podría emplear validación cruzada temporal $k$-fold para límites de confianza más ajustados.

\subsection{Robustez a Hiperparámetros Fijos}
Una preocupación natural con optimización por-escenario es si el desempeño TRSS es frágil a los valores específicos de $\alpha, \beta$ encontrados para cada nivel de degradación. Para abordar esto, evaluamos TRSS usando un único conjunto \emph{global} de hiperparámetros por dataset (la mediana de todos los óptimos por-escenario). En todos los 42 contextos experimentales, la configuración global exhibió un cambio de MAE promedio de $-0.7\%$ (es decir, negligentemente \emph{mejor} que sintonización por-escenario, debido a robustez de la mediana) y una pérdida SNR promedio de solo $0.46$~dB. Estos resultados demuestran que TRSS es inherentemente robusto a la elección de parámetros de regularización y no requiere re-optimización para cada patrón de fallo específico---una propiedad crítica para despliegue BCI en el mundo real donde la degradación es impredecible.

\subsection{Observaciones Concluyentes sobre Benchmarking Justo}
Históricamente, nuevos algoritmos son frecuentemente comparados contra líneas base configuradas por defecto, inflando mejoras percibidas. Al someter todos los métodos---incluyendo splines clásicos---a búsqueda rigurosa de hiperparámetros específica-por-dataset, este estudio establece un estándar de evaluación justo y exigente. El hecho de que TRSS y regularización de Tikhonov mantengan márgenes estadísticamente significantes (Bonferroni-corregido $p < 0.001$) sobre líneas base optimizadas valida la premisa fundamental: integrar topología de grafo espacial con priors de suavidad temporal constituye un framework estructuralmente superior para recuperación de señal EEG.
